%! Rational Functions and Holes — Unit 3, Lesson 3.4 — Lesson Plan
\documentclass[10pt]{article}
\usepackage{precalculus-article}
\usepackage{precalculus-boxes}
\usepackage{graphicx}
\graphicspath{{images/}}
\newcommand{\TallMath}[1]{$\displaystyle #1\rule[-1.4em]{0pt}{3.2em}$}
\newcommand{\CourseName}{Precalculus}
\newcommand{\MeetingLength}{60 minutes}
\newcommand{\UnitNumberName}{Unit 3: Rational Functions \quad}
\newcommand{\LessonNumberName}{Lesson 3.4: Rational Functions and Holes}

\begin{document}
\begin{center}
    {\Large\bfseries \CourseName \\
    {\small\bfseries \UnitNumberName \LessonNumberName}}
\end{center}

\begin{tcolorbox}[colback=sky, colframe=navy, boxrule=0.9pt, arc=2mm,
    left=2mm, right=2mm, top=1.5mm, bottom=1.5mm]
    \textbf{Primary Objective:} Students will identify holes in rational function graphs by
    locating shared zeros of the numerator and denominator where numerator multiplicity is
    greater than or equal to denominator multiplicity, simplify to find the hole's coordinates,
    and express the hole location using limit notation.
    \hfill\textit{\small Skill: MP 3.C}
\end{tcolorbox}

\renewcommand{\arraystretch}{1.6}
\begin{skillbox}[Priority Ideas \& Skills]{goldbox}
\begin{tabularx}{\linewidth}{@{}X X@{}}
\textbf{AP Skills} &
\textbf{Key Understandings} \\
\midrule
\textbf{MP 3.C} --- Support conclusions or choices with a logical rationale or appropriate data. &
\begin{itemize}[leftmargin=1.2em,itemsep=2pt,topsep=0pt]
  \item A hole occurs at $x = c$ when the numerator and denominator share a zero at $c$ with
        numerator multiplicity $m \geq$ denominator multiplicity $n$. (EK 1.10.A.1)
  \item The $y$-coordinate of the hole is $L = \lim_{x \to c} r(x)$, found by substituting
        $x = c$ into the simplified form. (EK 1.10.A.2)
  \item Both one-sided limits equal $L$: $\lim_{x \to c^-} r(x) = \lim_{x \to c^+} r(x) = L$.
\end{itemize} \\
\end{tabularx}
\end{skillbox}

\begin{skillbox}[{Vocabulary, Concepts, \& Theorems}]{greenbox}
\renewcommand{\arraystretch}{1.4}
\begin{tabularx}{\linewidth}{@{} >{\leavevmode\bfseries\color{navy}}p{0.27\linewidth} X @{}}
Hole in a graph         & A missing point $(c, L)$ where $r(c)$ is undefined but $\lim_{x \to c} r(x) = L$ exists. \\
Removable discontinuity & Another name for a hole; the discontinuity is ``removable'' by redefining the function at one point. \\
Common factor           & A polynomial factor shared by numerator and denominator; canceling it reveals a hole. \\
Multiplicity            & The number of times a given zero appears as a factor. \\
Limit notation          & $\lim_{x \to c} r(x) = L$: the outputs of $r$ approach $L$ as inputs approach $c$. \\
\end{tabularx}
\end{skillbox}

\begin{skillbox}[Activate Prior Knowledge \& Spiral Review]{sky}
    \begin{tabularx}{\linewidth}{@{}X@{}}
        Spiral review (warm-up) covers: identifying vertical asymptotes from 1.9 (shared zero,
        numerator nonzero or denom mult $>$ num mult); interpreting one-sided limit behavior at a
        VA; factoring a quadratic (bridge to factoring numerators and denominators for holes).
    \end{tabularx}
\end{skillbox}

\begin{skillbox}[Hook]{sky}
Suppose your GPS loses signal for exactly one moment --- at $t = 5$ seconds --- but your
position function is still defined everywhere else. The output \textit{at} $t = 5$ is missing,
but we can figure out where it ``should'' be by looking at what the function approaches from
both sides. This missing-but-predictable point is exactly what a hole is in a rational function.
\end{skillbox}

\begin{skillbox}[Lesson]{sky}
\begin{multicols}{2}
\textbf{Part 1 (10 min): Holes vs.\ Vertical Asymptotes.}

Given a shared zero $x = a$ with numerator multiplicity $m$ and denominator multiplicity $n$:
\begin{itemize}[leftmargin=1.2em,itemsep=2pt,topsep=2pt]
  \item $m \geq n$ $\Rightarrow$ \textbf{HOLE} at $x = a$
  \item $m < n$ $\Rightarrow$ \textbf{Vertical Asymptote} at $x = a$
  \item Zero only in denominator $\Rightarrow$ Vertical Asymptote
\end{itemize}

Connect to 1.9: VAs arise when denom is zero and numerator is nonzero \textit{or} when
$m < n$. Today's lesson is the complementary case ($m \geq n$).

\columnbreak

\textbf{Part 2 (15 min): Finding Hole Coordinates.}

Four-step procedure:
\begin{enumerate}[leftmargin=1.2em,itemsep=2pt,topsep=2pt]
  \item Factor numerator and denominator.
  \item Identify shared zeros and their multiplicities.
  \item Cancel the common factor.
  \item Substitute $x = c$ into the simplified form to find $L$.
\end{enumerate}

\textit{Example 1:} $r(x) = \frac{x^2-4}{x-2} = \frac{(x-2)(x+2)}{x-2}$.
Cancel $(x-2)$. Simplified: $x+2$. Hole at $(2, 4)$.

\textit{Example 2:} $r(x) = \frac{x^2-1}{(x-1)(x+2)}$. Cancel $(x-1)$.
Simplified: $\frac{x+1}{x+2}$. Hole at $\!\left(1, \frac{2}{3}\right)$. VA at $x=-2$.

\textbf{Part 3 (5 min): Limit Notation.}

If the hole is at $(c, L)$, then $\lim_{x \to c} r(x) = L$ with
$\lim_{x \to c^-} r(x) = \lim_{x \to c^+} r(x) = L$.
A hole is a \textit{removable} discontinuity: $r(c)$ undefined, but limit exists.
\end{multicols}
\end{skillbox}

\begin{skillbox}[Active Monitoring]{redbox}
\begin{itemize}[leftmargin=1.2em,itemsep=2pt,topsep=0pt]
  \item Cold-call after Part 1: ``Why does $m \geq n$ produce a hole instead of an asymptote?''
  \item Watch for students who cancel factors but forget to check the \textit{simplified} form
        for remaining VAs (e.g., Example 2 still has VA at $x = -2$).
  \item Confirm students substitute into the simplified expression, not the original.
\end{itemize}
\end{skillbox}

\begin{skillbox}[Group Work \& Differentiation]{redbox}
\begin{itemize}[leftmargin=1.2em,itemsep=2pt,topsep=0pt]
  \item \textbf{Tier R (Remediate):} Numerator/denominator pre-factored; classify hole vs.\ VA; fill table.
  \item \textbf{Tier A (Apply):} Factor independently; apply multiplicity rule; find coordinates $(c, L)$.
  \item \textbf{Tier E (Extend):} Write limit notation; explain reasoning; construct a rational function
        with specified holes and VAs.
\end{itemize}
\end{skillbox}

\begin{skillbox}[Individual Work \& Assessment]{redbox}
Exit ticket (2 questions, 1 page): factor $r(x) = \frac{x^2-x-2}{x^2-4}$, identify shared zeros,
classify each as hole or VA, find hole coordinates, and write the limit statement.
Collect and sort by Tier R/A/E for formative data before Lesson 1.11.
\end{skillbox}

\begin{skillbox}[Reinforcement \& Extension]{goldbox}
Homework: 5 problems covering identification of holes and VAs, multiplicity reasoning,
writing limit notation in words, a higher-multiplicity case ($x^3 - x^2$ over $x^2 - x$),
and an AP-style MC problem.
Preview of Lesson 1.11: equivalent representations of rational functions --- connecting
factored, standard, and graphical forms.
\end{skillbox}

\end{document}
